\documentclass{article}
\usepackage[margin=2cm]{geometry}
\usepackage{amsmath}
\usepackage{graphicx}
\usepackage{booktabs}
\usepackage[utf8]{inputenc}
\usepackage{ragged2e}
\setlength{\emergencystretch}{3em}
\begin{document}
\sloppy

\subsection*{0.1Mass-deformed USp(2N) gauge theory}


As discussed previously, we now give a mass to a single hypermultiplet in the fundamental
representation. This amounts to making a shift  ->  + m in the relevant functional
determinant. The result of this shift may be accounted for in by writing


\begin{equation}
\begin{array} { c c } { F ( \lambda _ { i }, m ) = \sum _ { i \neq j } \big [ F _ { V } ( \lambda _ { i } - \lambda _ { j } ) + F _ { V } ( \lambda _ { i } + \lambda _ { j } ) + F _ { H } ( \lambda _ { i } - \lambda _ { j } ) + F _ { H } ( \lambda _ { i } + \lambda _ { j } ) \big ] } & { } \\ { \qquad + ( N _ { I } - 1 ) F _ { H } ( \lambda _ { i } + m ) + F _ {}}\end{array}
\end{equation}


As before, we assume that A, = Nx for  > 0 and introduce a density p() satisfying .
Using the expansions , we find the analog of to be


\begin{equation}
\begin{array} { r } { F ( \mu ) \approx - \frac { 9 \pi } { 8 } N ^ { 2 + \alpha } \int d x d y ~ \rho ( x ) \rho ( y ) \left ( | x - y | + | x + y | \right ) + \frac { \pi } { 3 } ( 9 - N _ { f } ) N ^ { 1 + 3 \alpha } \int d z ~ \rho ( x ) ~ | x | ^ { 3 } } \\ { - \frac { \pi } { 6 } N ^ { 1 + 3 \alpha } \int d x ~ \rho ( x ) \left [ | x + \mu | ^ { 3 } + | x - \mu | ^ { 3 } \right ] ~ ( 2}\end{array}
\end{equation}


where for convenience we have defined  = m/N. As in the undeformed case, there is a
non-trivial saddle point only when  = 1/2. A normalized density function which extremizes
the free energy is


\begin{equation}
\begin{array} { r } { \rho ( x ) = \frac { 1 } { \left \{ 8 - N _ { f } \right \} z _ { a } ^ { 2 } - \mu ^ { 2 } } \left \{ 2 ( 9 - N _ { f } ) | x | - | x + \mu | - | x - \mu | \right \} \qquad x _ { s } = \sqrt { \frac { 9 + 2 \mu ^ { 2 } } { 2 ( 8 - N _ { f } ) } } \qquad }\end{array}
\end{equation}


with p() having support only on the interval x E [0, x*]. Inserting this result back into then
gives our final result,


\begin{equation}
\begin{array} { c c } { F ( \mu ) = \frac { \pi } { 1 3 5 } \left ( ( N _ { f } - 1 ) | \mu | ^ { 5 } - \sqrt { \frac { 2 } { 8 - N _ { f } } } ~ ( 9 + 2 \mu ^ { 2 } ) ^ { 5 / 2 } \right ) N ^ { 5 / 2 } \qquad } & { ( 4 ) } \\ {. } \end{array}
\end{equation}


We may check that when  = 0, we reobtain the result of the undeformed case . With this.
result and G6 given by we may now try to compare G6(F() - F(0)) to the same result
calculated holographically in Figure . Importantly, since  scales as N-1/2, we see that in
the large N limit the first term of is subleading and may be neglected. Thus to leading
order in N, the combination G6F() is in fact independent of Ng. Since comparison with.
the holographic result requires taking the large N limit, our supergravity solutions will be
unable to capture information about the precise flavor content of the SCFT dual. This agrees.
the n = 1 solutions we are considering can be consistently embedded into theories with any.

\end{document}
